\documentclass[10pt,a4paper]{article}
\usepackage[utf8]{inputenc}
\usepackage[T1]{fontenc}
\usepackage{helvet}
\renewcommand{\familydefault}{\sfdefault}
\usepackage[top=0.38in,bottom=0.38in,left=0.55in,right=0.55in]{geometry}
\usepackage{titlesec}
\usepackage{enumitem}
\usepackage{xcolor}
\usepackage{microtype}
\usepackage{ragged2e}
\usepackage{tikz}
\usepackage{graphicx}

\definecolor{primary}{RGB}{15,23,42}
\definecolor{accent}{RGB}{0,86,179}
\definecolor{sectionrule}{RGB}{132,132,132}
\definecolor{subtext}{RGB}{75,85,100}

\pagestyle{empty}
\setlength{\parindent}{0pt}
\setlength{\parskip}{0pt}
\setlength{\emergencystretch}{1.5em}
\raggedbottom

% Section formatting
\titleformat{\section}
  {\normalsize\bfseries\color{primary}}{}{0pt}{\MakeUppercase}
  [\vspace{-0.55em}\textcolor{sectionrule}{\rule{\textwidth}{0.5pt}}\vspace{0.06em}]
\titlespacing*{\section}{0pt}{0.32em}{0.12em}

% Header with circular photo
\newcommand{\cvheaderwithphoto}[6]{%
  \noindent
  \begin{minipage}[c]{0.80\textwidth}
    {\fontsize{20}{23}\selectfont\bfseries\color{primary} #1}\\[0.15em]
    {\small\bfseries\color{accent} #2}\\[0.20em]
    {\footnotesize #3}\\[0.10em]
    {\footnotesize #4}\\[0.10em]
    {\footnotesize #5}
  \end{minipage}%
  \hfill
  \begin{minipage}[c]{0.18\textwidth}
    \raggedleft
    \begin{tikzpicture}
      \clip (0,0) circle (1.18cm);
      \node[inner sep=0pt] at (0.45cm, -0.55cm) {\includegraphics[height=3.4cm,keepaspectratio]{#6}};
      \draw[draw=sectionrule!50, line width=0.8pt] (0,0) circle (1.18cm);
    \end{tikzpicture}
  \end{minipage}
  \vspace{0.15em}
}

% Experience item macro
\newcommand{\experienceitem}[5]{%
  \noindent\textbf{#1}\hfill\textbf{#2}\\[0.02em]
  {\footnotesize\textit{\color{accent} #3}}\hfill{\footnotesize\textbf{#4}}\\[0.04em]
  \if\relax\detokenize{#5}\relax\else{\scriptsize\color{subtext}\textbf{Technologien / Stack:} #5}\\[0.08em]\fi
}

% Education item macro
\newcommand{\educationitem}[4]{%
  \noindent\textbf{#1}\hfill\textbf{#2}\\[0.02em]
  {\footnotesize\textit{\color{accent} #3}}\hfill{\footnotesize\textbf{#4}}\\[0.04em]
}

% Project item macro
\newcommand{\projectitem}[4]{%
  \noindent\textbf{#1}\hfill{\footnotesize\textbf{#2}}\\[0.02em]
  {\scriptsize\color{subtext}\textbf{Stack:} #3}\hfill{\scriptsize\textbf{#4}}\\[0.06em]
}

% List formatting without full stops
\newcommand{\cvliststart}{%
  \begin{itemize}[leftmargin=1.15em,labelsep=0.38em,topsep=0.04em,itemsep=0.04em,parsep=0pt,partopsep=0pt]
}
\newcommand{\cvlistend}{\end{itemize}\vspace{0.14em}}

\newcommand{\competency}[1]{{\textbullet\enspace #1\enspace}}
\newcommand{\skillrow}[2]{\textbf{#1:}\enspace #2\\[0.10em]}

\begin{document}

% ================= PAGE 1 =================
\cvheaderwithphoto
  {Pradeep Gandu}
  {Cloud Data Engineer \enspace\textbar\enspace Automotive Software \enspace\textbar\enspace KI / Machine Learning}
  {Bernsdorfer Str. 115, 09126 Chemnitz \enspace\textbar\enspace +49 176 59677415 \enspace\textbar\enspace gandupradeep2026@gmail.com}
  {LinkedIn: linkedin.com/in/pradeep-gandu \enspace\textbar\enspace GitHub: github.com/gandupradeep2026 \enspace\textbar\enspace Portfolio: gandupradeep2026.github.io}
  {Deutsch: C1 (UNIcert III / DSH) \enspace\textbar\enspace Englisch: Fließend (C1) \enspace\textbar\enspace Status: § 16b AufenthG (Vollzeit ab 01/2027)}
  {photo.jpeg}

\section{Kurzprofil}
{\small
Ingenieur für Cloud-Datenplattformen, Künstliche Intelligenz und Automotive Software Engineering mit fundierter Praxiserfahrung in skalierbaren Streaming- und Batch-Pipelines (GCP BigQuery, Pub/Sub, Apache Beam, Dataflow) sowie generativer KI und modularen RAG-Architekturen (FastAPI, MCP, Ollama bei der DiPP GmbH). Masterand an der Technischen Universität Chemnitz mit verhandlungssicheren C1-Deutschkenntnissen, bundesweiter Mobilität und sofortiger Verfügbarkeit (ab sofort als Werkstudent / Teilzeit, ab 01/2027 in Vollzeit).
}

\section{Kernkompetenzen}
\vspace{0.06em}
{\small
\competency{Cloud Data Engineering (GCP / AWS)}\competency{Apache Beam / PySpark / Spark SQL}\competency{Fahrzeugtelemetrie \& CAN-Bus (ISO 14229)}\\
\competency{CARLA AD Simulator \& Scenario Runner}\competency{Python / FastAPI / REST-APIs}\competency{Generative KI / RAG / MCP / Ollama}\\
\competency{BigQuery / Cloud Storage / Cloud Composer}\competency{CI/CD GitHub Actions / Docker / Linux}
}
\vspace{0.10em}

\section{Berufserfahrung}

\experienceitem{Werkstudent — KI-Softwarelösungen}{DiPP GmbH}{Chemnitz}{01/2025 – 07/2025}{Generative KI, FastAPI, RAG, Model Context Protocol (MCP), Ollama, IONOS Cloud, Docker, Git}
\cvliststart
  \item Konzeption und Implementierung generativer KI- und LLM-Services für Enterprise-CRM-Systeme via Python und FastAPI
  \item Aufbau modularer RAG-Pipelines mit Ingestion, Text-Chunking, Vektor-Embeddings und semantischem Retrieval
  \item Standardisierung von Systemkontexten und Toolschnittstellen über das Model Context Protocol (MCP) für sichere lokale Inferenz
  \item Durchführung von Integrationsprüfungen, Fehleranalysen und systematischer Evaluation der Modellantworten
\cvlistend

\experienceitem{Associate Language Specialist (Content Engineering \& NLP)}{GlobalLogic Technologies}{Hyderabad, Indien}{10/2020 – 12/2021}{Deutsche Handelsdaten, NLP / Textverarbeitung, Machine Learning, scikit-learn, PyTorch, Python}
\cvliststart
  \item Bereinigung, Kuration und semantische Validierung deutscher Einzelhandelsdaten für automatisierte ML-Pipelines
  \item Training, Hyperparameter-Optimierung und Evaluation überwachter Klassifikationsmodelle mit Python und scikit-learn
  \item Durchführung systematischer Qualitätsprüfungen zur Sicherstellung hoher Annotationskonsistenz nach Kundenrichtlinien
\cvlistend

\experienceitem{Senior Process Executive (ML-Qualitätsprüfung Meta / Facebook)}{Cognizant}{Hyderabad, Indien}{09/2017 – 05/2019}{Facebook Marketplace, Human-in-the-Loop ML, Datenqualität, SQL, Python, Ursachenanalyse}
\cvliststart
  \item Durchführung von Human-in-the-Loop ML-Evaluationen und Metadatenvalidierungen für Facebook Marketplace
  \item Steigerung der Erkennungs- und Klassifikationsgenauigkeit von ca. 64\,\% auf 76\,\% durch datengestützte Feedbackschleifen
  \item Automatisierte Datenqualitätsaudits und strukturierte Ursachenanalysen bei Modellabweichungen mittels SQL und Python
\cvlistend

\experienceitem{Spezialist für CRM-Operations (Qualitätssicherung für Google)}{Concentrix Daksh Services}{Hyderabad, Indien}{03/2016 – 05/2017}{Google Ads, Workflow-Automatisierung, Prozessvalidierung, Fehleranalyse, SQL}
\cvliststart
  \item Prozessvalidierung, Fehlerauswertung und Datenprüfungen für automatisierte Google-Ads-Workflows
  \item Steigerung der operativen Prozess-Erfolgsquote von ca. 72\,\% auf 79\,\% durch Optimierung der Prüfschritte
\cvlistend

\newpage
% ================= PAGE 2 =================

\section{Ausbildung}

\educationitem{M.Sc. in Automotive Software Engineering}{Technische Universität Chemnitz}{}{10/2022 – voraussichtlich 12/2026}
\cvliststart
  \item 90 / 120 ECTS absolviert; Masterarbeit in Finalisierung: Generierung und Annotation szenariozentrierter Datensätze für Umgebungsvorhersagen beim autonomen Fahren (Python, CARLA Simulator 0.9.15, OpenDRIVE, Parquet)
  \item \textbf{Studienschwerpunkte:} Verteilte Systeme, Software-Architekturen, Fahrerassistenzsysteme, Simulationstechnik
\cvlistend

\educationitem{M.Tech. in Eingebettete Systeme (Embedded Systems)}{MLR Institute of Technology}{Hyderabad, Indien}{09/2017 – 09/2019}\\[-0.05em]
{\footnotesize First Class with Distinction (CGPA: 7,33 / 90 Credits) \enspace\textbar\enspace \textbf{Schwerpunkte:} RTOS, Mikrocontroller, Systementwurf}\\[0.20em]

\educationitem{B.Tech. in Elektro- und Informationstechnik}{Jawaharlal Nehru Technological University}{Hyderabad, Indien}{08/2009 – 11/2013}\\[-0.05em]
{\footnotesize First Class (Gesamtnote: 62,67\,\%) \enspace\textbar\enspace \textbf{Schwerpunkte:} Regelungstechnik, Elektronik, Signalverarbeitung}\\[0.25em]

\section{Ausgewählte Praxisprojekte}

\projectitem{Automotive-Telemetrie-Streaming-Plattform (GCP)}{2026}{GCP Pub/Sub, Apache Beam, Dataflow, BigQuery, Airflow, Python}{GitHub}
\cvliststart
  \item Skalierbare Event-Streaming-Pipeline für Fahrzeugtelemetriedaten mit zustandsbehafteter Deduplizierung via Apache Beam
  \item Partitionierte, kostenoptimierte Speicherung in Google BigQuery sowie automatisierte Orchestrierung mit Cloud Composer
\cvlistend

\projectitem{AutoSense AI: Vorausschauende Fahrzeugdiagnose \& Agent}{2026}{FastAPI, IsolationForest, RAG, Ollama LLM, MCP, React / Vite}{GitHub}
\cvliststart
  \item Anomalieerkennung an 39.000+ realen OBD-II-Datensätzen mit 5-aus-7-Filter zur Fehlalarmreduktion um 93,4\,\%
  \item Interaktiver KI-Diagnose-Agent mit RAG-Abfrage technischer Reparaturleitfäden und Model Context Protocol (MCP)
\cvlistend

\projectitem{TU Chemnitz Badminton Community Portal}{2026}{React / Vite, Node.js / Express, SQLite, Tailwind, Cloudflare, Oracle VM}{Live-App}
\cvliststart
  \item Entkoppelte Full-Stack-Webanwendung auf GitHub Pages mit Node.js/Express REST-Backend via Cloudflare-Tunnel
  \item Sicheres Turnieranmeldesystem, Trainervermittlung und gehärtete Authentifizierung mit kryptografischen 6-Digit-OTPs
\cvlistend

\section{Fähigkeiten \& Qualifikationen}
{\small
\skillrow{Cloud \& Big Data}{Google Cloud Platform (BigQuery, Pub/Sub, Dataflow, Dataproc, GCS, IAM), AWS (EC2, S3), Apache Beam, PySpark, Spark SQL, Apache Airflow, Parquet}
\skillrow{Software \& KI}{Python (FastAPI, Scikit-learn, PyTorch), SQL (Experte), RAG-Architekturen, Ollama, MCP, REST-APIs, Docker, Docker Compose, Git / GitHub Actions CI/CD}
\skillrow{Automotive \& IT}{CARLA AD Simulator, OpenDRIVE, 3D-Trajektorien, CAN-Bus, ISO 14229 (UDS), Linux (Ubuntu/Debian), Windows Server, Active Directory}
\skillrow{Zertifizierungen}{Google Cloud Certified Data Engineer \enspace\textbullet\enspace UNIcert III C1 Deutsch (TU Chemnitz) \enspace\textbullet\enspace Cisco Python \enspace\textbullet\enspace HackerRank SQL}
}

\section{Sprachkenntnisse}
{\small
\textbf{Deutsch:} C1 (UNIcert III / DSH, verhandlungssicher) \enspace\textbar\enspace \textbf{Englisch:} C1 (Fließend in Wort und Schrift) \enspace\textbar\enspace \textbf{Telugu:} Muttersprache \enspace\textbar\enspace \textbf{Hindi:} Fließend
}

\section{Engagement \& Interessen}
{\small
\textbf{Badminton:} Aktiver Spieler im Hochschulsport der TU Chemnitz sowie Initiator und Entwickler des universitären Badminton-Community-Portals\\
\textbf{Interkultureller Austausch:} Aktive Unterstützung und Vernetzung internationaler Masterstudierender an der TU Chemnitz
}

\vspace{0.6em}
\noindent
\begin{minipage}[b]{0.48\textwidth}
  Chemnitz, den 14. September 2026\\[1.0em]
  \rule{5.5cm}{0.4pt}\\
  {\footnotesize Ort, Datum}
\end{minipage}%
\hfill
\begin{minipage}[b]{0.48\textwidth}
  \raggedleft
  \includegraphics[height=0.90cm,keepaspectratio]{Sign.png}\\[0.05em]
  \rule{5.5cm}{0.4pt}\\
  {\footnotesize Pradeep Gandu (Unterschrift)}
\end{minipage}

\end{document}
